
\subsection{Lay Activity According to Ability and Gift}

When anyone in the United Church, which has after all set itself the task of “fostering Christian lay activity,” ventures to give a little encouragement toward the exercise of this work, he must put up with both mockery and abuse from the official organ of the church body. It is “foolishness” and “enthusiasm” and “ecclesiastical anarchism.” — Yes, it is now possible to understand how our church assemblies have managed to arrange matters so that, as a rule, the Christian laity sit silent.

Since the matter has now come up for serious discussion, however, we shall try to speak straight from the heart concerning a matter serious enough for the Christians among us, however much it may be turned into a cudgel for church politicians.

Life and Spirit are the chief thing in all Christianity, and therefore also in all ecclesiastical Lutheranism. Life and Spirit are the chief thing in all Christian undertakings; therefore also in lay activity. A Christian activity that is merely learned and practiced, but lacks the Holy Spirit, is neither pleasing to God nor a blessing to the congregation.

According to Scripture, Christian lay activity does not have its root, ground, and source in certain particular gifts, whether one calls them extraordinary or ordinary spiritual gifts. According to Scripture, the source and fountain of Christian lay activity lie in the living experience of salvation itself, in the impartation of the Spirit itself in regeneration, in the very passage from darkness into God’s marvelous light. Whoever does not know this knows neither Christ nor the power of God as it is experienced in personal salvation, or in the passage from death to life.

According to Scripture, on the other hand, lay activity receives its limitation and all its more particular determination from the particular spiritual gift or gifts that each believer has received.

For the calling to witness, which is the fountain and source of lay activity, is universal and belongs to and rests upon all living Christians; but the manner in which I am to exercise my calling to witness depends upon my capacity of utterance and my spiritual gift. All living Christians must exercise lay activity; but not all in the same manner; that is determined and limited by ability and gift.

This is according to Scripture. Already in the Old Testament we find it expressed clearly enough, time after time, that whoever is saved by the Lord must fulfill his vows to the Lord and honor his name by confessing his saving wonders.

Over the people of Israel in their childhood Moses uttered the wish: “Would that all the Lord’s people were prophets, and that the Lord would put his Spirit upon them!” Num. 11:29. The old man of God was jealous of no one and cherished no petty fear of lay activity. If only there could be enough of it — by the Spirit of the Lord!

Therefore this also stands as the Lord’s faithful promise over his people in the last days: “And it shall come to pass afterward that I will pour out my Spirit upon all flesh, and your sons and your daughters shall prophesy,” Joel 3:1. Moses’ wish shall be fulfilled, and a people of God shall be created in which all true members shall partake of the Spirit and of prophecy.

Thus David also speaks of the proper thanksgiving for salvation, which therefore rests upon everyone who is saved — each according to ability and gift:

“I will declare your name to my brothers; in the midst of the congregation I will praise you.” Ps. 22:23.

“From you comes my praise in the great congregation; I will fulfill my vows before those who fear him.” Ps. 22:26.

“Rescue my soul from their destructions, my only one from the young lions. I will give thanks to you in the great congregation; I will praise you among many people.” Ps. 35:17–18.

“I proclaimed righteousness in the great congregation; see, I did not close my lips; Lord, you know it. I did not hide your righteousness within my heart; I declared your faithfulness and your salvation; I did not conceal your mercy and your truth from the great congregation.” Ps. 40:10–11.

“Come and hear, all you who fear God, and I will tell what he has done for my soul.” Ps. 66:16.

“Hallelujah, I will praise the Lord with my whole heart, in the council of the upright and in the congregation.” Ps. 111:1.

“What shall I render to the Lord for all his benefits toward me? I will lift the cup of salvation and call upon the name of the Lord. I will fulfill my vows to the Lord, yes, before all his people.” Ps. 116:12–14.

“I believed, therefore I spoke.” Ps. 116:10.

These passages from the Old Testament, and many more with them, show quite irresistibly that salvation obligates one to audible confession in the congregation of God. The Lord’s great works, which the individual believer experiences, lay vow and obligation and responsibility upon him, and he must come forth with an offering of thanksgiving and with a testimony of praise in the courts of the Lord and in the assembly of his congregation.

Should there then be less duty to witness, less praise, less right, and less freedom in the New Testament? Can we now be saved from death and carry salvation in silent hearts and with closed lips? Or is praise not due to the Lord, and do not many fellow travelers need to hear the testimony? Shall we fall behind those whom the Lord saved under the old covenant?

Let us look at the testimony of the New Testament in this matter!

The Savior appoints his congregation to be the salt of the earth and the light of the world, a city on the mountain that cannot be hidden, made manifest through the testimony of the truth and the works of love, Matt. 5:13–16. Is there anyone who believes this can happen without the participation and cooperation of all believers?

The Savior says: “Therefore everyone who will confess me before men, him will I also confess before my Father who is in heaven.” Matt. 10:32. Thus men and women, all who believe, are equally obligated — under peril of the loss of salvation — to confess Jesus before men.

But someone answers: Yes, all must witness by their lives, by their beautiful conduct, by their good works; not by word and speech. But then Paul comes and says: “For if you confess with your mouth that Jesus is Lord, and believe in your heart that God raised him from the dead, then you shall be saved. For with the heart one believes unto righteousness, but with the mouth one confesses unto salvation,” Rom. 10:9–10.

And Peter makes all believers proclaimers and preachers when he says: “But you are a chosen race, a royal priesthood, a holy people, a people for his own possession, that you may proclaim the excellencies of him who called you out of darkness into his marvelous light.” 1 Pet. 2:9.

Salvation by faith therefore creates “the fruit of lips,” as the Epistle to the Hebrews also says: “Let us then always through him offer to God the sacrifice of praise, that is, the fruit of lips that confess his name.” Heb. 13:15.

And that proclamation has no other ground or source in any right and true preacher, whether he be lay or learned, apostle or member of the congregation, Paul expresses clearly when he writes concerning himself to the Corinthians:

“But since we have the same Spirit of faith according to what is written: I believed, therefore I spoke — we also believe, therefore we also speak,” 2 Cor. 4:13.

Paul speaks, driven by the same Spirit of faith that all believers also have. When he became a believer, he became a witness of Christ. Then there was poured out within him that love of Christ which afterward became the compelling power in the work of his life.

Therefore Paul also desires to have all believers taking part in the work and testimony of the congregation; and his “liturgy,” or form of worship, is entirely simple and free:

“What then is to be done, brothers? When you come together, each of you has a psalm, he has a teaching, he has a tongue, he has an interpretation; let all things be done for edification.” 1 Cor. 14:26.

There is not much stiffness in Paul’s “liturgy.” Nor is lay activity excluded; for there is in fact nothing else here than what we call lay activity. Therefore he also says: “For you can all prophesy, one by one, so that all may learn and all may be exhorted.” 1 Cor. 14:31.

But that lay activity is to be exercised according to ability and gift, Paul expresses with striking clarity:

“There are differences of spiritual gifts, but the Spirit is the same; and there are differences of ministries, but the Lord is the same; and there are differences of mighty workings, but God is the same, who works all things mightily in all. But the manifestation of the Spirit is given to each one for what is profitable. For to one is given a word of wisdom through the Spirit, but to another a word of knowledge through the same Spirit, and to another faith through the same Spirit, and to another gifts of healing through the same Spirit, and to another the working of miracles, and to another prophecy, and to another the discerning of spirits, and to another various kinds of tongues, and to another the interpretation of tongues. But all these things the one and same Spirit works mightily, distributing to each one individually as he wills.” 1 Cor. 12:4–11.

And in another place:

“But since we have different spiritual gifts according to the grace that has been given us, whether we have the gift of prophecy, let us use it in accordance with faith; or if we have a ministry, let us attend to the ministry; or if someone is a teacher, to the work of teaching; or if someone exhorts, to the exhortation! The one who gives, let him do it in sincerity; the one who leads, with zeal; the one who shows mercy, let him do it with gladness!” Rom. 12:6–8.

See how lively and stirred, how free and spiritual Paul wished the congregation to be! What movement! What life! What unfolding of abilities and gifts, what freedom!

This also follows of necessity from the biblical fundamental view of the matter, namely, that whoever is saved is, by salvation itself, called to be a witness of the Lord and a fisher of men. For if something is a duty and a demand of God, then surely the rule is this: From the one to whom little is given, little shall be required; but from the one to whom much has been entrusted, more shall be demanded.

Or is one person to love less than he receives Spirit and grace to do, and another more? If love is God’s gift and God’s demand, then surely the rule is simple enough: the more, the better.

So also with the duty to witness. We have merely selected a few passages of Scripture to show how the matter stands. Some time ago we believed it quite superfluous among Christians to cite any Scripture at all in order to demonstrate so simple a matter. Norwegian believers are, after all, familiar with all this. But our experience of the opposition to free Christian lay activity has convinced us that, unfortunately, it is still necessary to lay the foundation in this matter.

God willing, we shall later show how the experience of God’s children agrees wholly and entirely with Scripture on this point.

---

When the attacks upon free Christian lay activity assumed so acute a form that it was called enthusiasm and anarchism, I found it necessary in the last issue of “Folkebladet” to demonstrate the broad and solid biblical ground upon which lay activity rests. It stands and falls with the Christian duty to witness. That it assumes different forms and appears more or less publicly, in small and large assemblies, in home, school, and church, is shown by both the Bible and church history. But its fountain and source are in the Christians’ calling to witness.

We come now to Christian experience. For there is no duty that can be laid upon the believer unless it is also wrought in him by the Holy Spirit. Therefore Christian experience will also cast light upon all Christian matters. Christian duty and Christian freedom agree with one another, and we unite them both when we speak of the Christian calling; for within it lie both what God requires of them and what he himself by the Spirit impels them to do.

The experience of believers in this matter is namely this: they themselves, having been saved by God’s love, are driven by the love that saved them to have others saved with them.

With all the inwardness of experience this has been described once for all in John 1:35–52, so that it would be almost wrong, and in any case superfluous, to relate how the same thing has repeated itself in everyone who has come to faith. John tells how those who found Jesus had to find, each one, his brother and his friend and say to them: We have found the Messiah; come and see!

Joy over salvation and living, overflowing love drive these believing men to find more people who need the same deliverance that they themselves have experienced. They have found the Lamb of God, who bears the sins of the world, and how many there are who need to find him! Therefore they must go out and bring the great, glad message first to those nearest them.

Come and see! That is the true experiential preaching, and it is and must be the fundamental note in lay preaching. But this is also precisely what every single believer both can and must say in his confession of Jesus.

For there are two things which, above all else, force themselves upon the consciousness of the one who has been saved; the one is the dreadful misery and terrible danger from which he has been saved; the other is the surpassing grace and blessed joy that he has found with Jesus. And both these experiences drive the saved person to find those who need the same salvation as he himself.

To see people who stand near us in life walking blindly upon the broad road to hell is utterly unbearable for believing people. We must warn them, whether it is called order or disorder. The sleeping church may take offense as much as it pleases.

And the knowledge of Christ’s love and grace is so surpassingly joyful that it is impossible not to let sinners hear of it. Believing people must confess, so that others may become as blessed as they themselves.

But how this calling to witness, which every believer bears within himself, is to be exercised to the greatest blessing depends upon the Spirit’s prompting, or upon ability and gift. Not all see equally far, and not all are driven far afield. Hans Nielsen Hauge had to go out to the furthest boundaries of Norway. And, God willing, another time we shall write a little about how lay activity expressed itself in his day.

For now, the chief thing is to take hold in earnest of all the sleeping Christians and congregations among us, in the hope that perhaps the Lord may grant that they be awakened.

\subsection{Lay Activity in Dead Congregations}

Since lay activity arises from the Christian calling to witness, which is given to all believers when they come to faith; and since this calling to witness has above all else the purpose of awakening those who sleep from the danger of death and condemnation before it is too late, it will easily be understood that true Christian lay activity has a particular need to break in with the awakening call precisely where the darkness is thickest and the sleep heaviest, and therefore the danger greatest.

Where lay activity does not have this character within it, it suffers from so serious a fundamental defect that it is a question whether there is much value in it at all. Therefore it is particularly necessary for those who travel about speaking the Word of God to examine themselves honestly and thoroughly as to whether what matters most to them is to bring the testimony to places where it is dark and cold and dead, or whether they prefer to seek out those congregations where believing friends receive them with joy and gifts.

But can lay preachers really enter the dead congregations, where people entrench themselves behind “right doctrine” and want no lay activity? Should not the layman who travels with the confession of God’s Word respect congregational resolutions and pastoral commands and keep far away from such congregations?

No, far from it. If Hans Nielsen Hauge had waited for permission from congregations and pastors, there would have been no lay activity. And it will not do to let people sleep in peace merely because they call themselves a congregation and keep a dead pastor to help them sleep more securely. Or when people are sleeping in a burning house, does the person who is awake have no right to cry “Fire!” until he has received a call to do so from those who are sleeping?

A multitude of congregations have matters arranged in such a way that they use the pastor and the worship service and their entire ecclesiastical apparatus as a defense and a screen against the Spirit of God and living Christianity. And perhaps no congregation is entirely free from this sin; for the snare is so subtle that it steals in upon even the most watchful in order to take them by surprise. Thus there are also many congregations that, by means of their very congregational organization and ecclesiastical order, defend themselves against lay activity. Many ask to this very day: Why then do you baptize, if you are not Christ, nor Elijah, nor the Prophet? They want order, ecclesiastical order; the voice crying in the wilderness must not disturb them or their order.

Must one bow before this order? Shall those who sleep be allowed to sleep in peace because they have such an order that no one is to have the right to disturb them? Impossible, if there is otherwise any right belonging to love. That people use God’s holy congregation as a shield against the Spirit’s awakening voice is a sin that by no means gives them a right to escape the preaching of repentance, if there is any way of getting it into their ears and hearts.

If, therefore, there is any meaning in a renewal of Christian lay activity among us, then above all let us see it attack the dead congregations, where there is more scorn to be gained than any other reward. If one hears of a congregation that closes its church to lay preachers, or of a pastor who regards Christian lay activity as sin, that is precisely the place where it is more important than elsewhere for the crying voice to be heard. Go in there, however loud the outcry and protest may become. “It was as though I saw the world lying submerged in evil,” Hauge says of his calling. Therefore he would go.

But here also it holds true that the activity becomes “according to ability and gift.” For no one can labor farther than his vision and his task extend. Hauge saw out over all Norway, and nothing could hold him back from breaking his way forward to the outermost boundaries of the land. He had to press forward everywhere, and his faith sustained him immovably until the prison wall shut itself between him and his field of labor. Others may perhaps see their task more narrowly: their house, their family, their congregation, their pastor, their neighbor, their Sunday school they recognize, at least, as their nearest field of labor. Not all Christians must travel about with the testimony. But often those who remain at home must help others to go out. And not least when the matter is one of pressing in with the Word where “ecclesiastical order” places such heavy obstacles in the way.

Where there is wilderness, there the crying voice shall be heard. Brothers, where is the wilderness? Or rather: Where is the wilderness not? 
